% =====================================================================
%  Presenter notes for "Online PAC Labeling from an LP Perspective"
%  Companion to slides.tex  (~15 minute talk)
%
%  Each section corresponds to one slide, in order. Bold blue cues are
%  the spoken throughline; the "Backup" paragraphs are deeper technical
%  detail to draw on. A final section anticipates audience questions.
%
%  Compile with:  pdflatex presenter_notes.tex
% =====================================================================
\documentclass[11pt]{article}
\usepackage[margin=1in]{geometry}
\usepackage{amsmath,amssymb,mathtools}
\usepackage{xcolor}
\usepackage{enumitem}

\newcommand{\eps}{\varepsilon}
\newcommand{\R}{\mathbb{R}}
\newcommand{\Pp}{\mathbb{P}}
\newcommand{\calA}{\mathcal{A}}
\newcommand{\calG}{\mathcal{G}}
\newcommand{\argmin}{\operatorname*{arg\,min}}
\newcommand{\hath}{\widehat h}
\newcommand{\hatL}{\widehat L}
\newcommand{\tildeY}{\widetilde Y}
\newcommand{\hatY}{\widehat Y}
\newcommand{\E}{\mathbb{E}}
\newcommand{\ind}[1]{\mathbf{1}\{#1\}}
\newcommand{\cue}[1]{\textcolor{blue!60!black}{\textbf{#1}}}
\newcommand{\timecue}[1]{\hfill{\small\itshape(\,$\approx$#1\,)}}
\newcommand{\backup}[1]{{\small\textcolor{black!72}{\textit{Backup --- } #1}}}

\renewcommand{\thesection}{\arabic{section}}
\renewcommand{\section}[1]{%
  \refstepcounter{section}%
  \par\vspace{0.6\baselineskip}%
  {\large\bfseries\color{blue!45!black}Slide \thesection. #1}\par
  \vspace{0.2\baselineskip}}
\newcommand{\qsection}[1]{%
  \par\vspace{0.7\baselineskip}%
  {\large\bfseries\color{blue!45!black}#1}\par
  \vspace{0.2\baselineskip}}
\setlength{\parskip}{0.45\baselineskip}
\setlength{\parindent}{0pt}

\title{\vspace{-1.2cm}Presenter Notes\\[2pt]\large Online PAC Labeling from a Linear Programming Perspective}
\author{Companion to \texttt{slides.tex} --- target $\approx$15 minutes}
\date{}

\begin{document}
\maketitle
\vspace{-0.8cm}

\noindent\textbf{Pacing guide.} Setup + LP (0:00--4:00), offline + results
(4:00--11:30), then the \emph{online segment} (Slides 12--17) as a self-contained
$\approx$5-minute block. Parenthetical times are \emph{cumulative} through Slide
11; for the online block they restart and are \emph{relative to the start of the
segment} (it is a separate speaker). If running long, compress the dual
derivation (Slide 5) and the proof sketches (Slides 9 and 15) and lean on the
figures --- those are the parts an audience can skip without losing the thread.

\noindent\textbf{The one-sentence thesis} (keep returning to it): \emph{a single
LP dual variable $\lambda$ --- the ``price of error'' --- unifies offline
certification and online auditing, and proxies only ever affect cost, never the
error guarantee.}

% --------------------------------------------------------------------
\section{Title}
\timecue{0:20}
\cue{Open simply.} ``Our project revisits PAC labeling --- cheaply labeling a
large dataset with a finite-sample error guarantee --- but through the lens of a
linear program. That single change of viewpoint buys us multi-source routing and
an online extension almost for free.'' Name the three authors, then move on.

\backup{``PAC'' = probably approximately correct: the error guarantee holds with
probability $1-\alpha$ over the calibration (offline) or audit (online)
randomness. Our starting point is PAC labeling of Cand\`es et al.\ (2025); our
contribution is the optimization-first framing, item/source-specific costs, and
the audited online extension.}

% --------------------------------------------------------------------
\section{The problem: labeling a big pool, cheaply}
\timecue{1:30}
\cue{Set the stakes.} We have a fixed pool of $T$ unlabeled items
$X_1,\dots,X_T$ with unknown true labels $Y_1,\dots,Y_T$, and we must release a
label for \emph{every} item. Two ways to get a label: $k$ cheap ML models (free
but fallible) or a human \textbf{expert} (treated as ground truth, but
expensive). We want to lean on the AI as much as possible while
\emph{certifying} that the average error of the released dataset is below a
user-chosen target $\eps$.

\cue{State the tension --- it organizes the whole talk.} More AI lowers cost;
more expert lowers error. Everything that follows is about navigating that
frontier \emph{with a guarantee}, not just a heuristic. The headline numbers to
foreshadow: at a 10\% error target we skip the expert on roughly 78\% of
ImageNet items.

\cue{Why ``LP''?} Because the trade-off is literally a constrained
cost-minimization: minimize expert queries subject to an error budget. Writing
it as an LP and taking the dual is what collapses a $T\times(k{+}1)$ decision
problem down to tuning one scalar.

% --------------------------------------------------------------------
\section{Setup and notation}
\timecue{3:00}
\cue{Walk the two boxes.} Action set $\calA=\{E,1,\dots,k\}$: defer to the
expert $E$, or accept one of $k$ cheap models that produce predictions
$\hatY_t^{(j)}$. Two item-wise primitives:
\[
   h_t(a)=\begin{cases}0,&a=E\\ \ell(Y_t,\hatY_t^{(a)}),&a\in[k]\end{cases}
   \qquad
   c_t(a)=\begin{cases}1,&a=E\\ 0,&a\in[k].\end{cases}
\]
\cue{Emphasize the asymmetry --- it drives every later simplification.} The
expert has zero loss but unit cost; cheap models have zero cost but possibly
positive loss. The labeling loss $\ell$ is bounded and normalized to $[0,1]$.
The released label is $\tildeY_t=Y_t$ when we defer, else $\hatY_t^{(a_t)}$.

\backup{The unit/zero cost split is the ``equal cheap-source cost'' regime
($c_t(j)=0$ for all $j$, $c_t(E)=1$). The \emph{LP itself} permits item- and
source-specific costs; we specialize to equal costs because that is what makes
$L_T(\lambda)$ monotone and what the data supports. Unequal costs come back in
the limitation slide and in the Q\&A.}

% --------------------------------------------------------------------
\section{The cost-aware LP}
\timecue{4:00}
\cue{Read the LP in words, not symbols.} Decision variable
$x_{t,a}\in[0,1]$ is the probability of taking action $a$ on item $t$. We
\emph{minimize average cost} $\frac1T\sum_{t,a}c_t(a)x_{t,a}$ subject to: (i) the
\emph{average error budget} $\frac1T\sum_{t,a}h_t(a)x_{t,a}\le\eps$; (ii) each
item's probabilities sum to one; (iii) nonnegativity.

\cue{Two things to stress.} First: there is exactly \emph{one} coupling
constraint --- the global error budget --- and that single coupling is what
makes the dual one-dimensional. Everything else is per-item (a simplex). Second:
this is the \emph{oracle} LP --- if the true losses $h_t(a)$ were known, its
solution \emph{is} the optimal labeling policy. The rest of the talk is about not
knowing $h_t(a)$ and plugging in proxies $\hath_t(a)$.

\backup{The LP is a continuous relaxation, but its optimum is attained at a
vertex, so the per-item solution is essentially deterministic (an item is
routed, deferred, or sits exactly at a threshold). This is why the eventual rule
is a hard threshold, not a randomized mixture.}

% --------------------------------------------------------------------
\section{The dual / price of error}
\timecue{5:00}
\cue{This is the conceptual heart --- slow down here.} Put a single multiplier
$\lambda\ge0$ on the error constraint. The reduced dual is the one-dimensional
\emph{concave} maximization
\[
   \max_{\lambda\ge0}\;D(\lambda):=-\lambda\eps
     +\frac1T\sum_{t=1}^{T}\min_{a\in\calA}\{c_t(a)+\lambda\,h_t(a)\}.
\]
The structural payoff: \emph{for fixed $\lambda$ the inner minimization
separates across items}, giving the fixed-price rule
$a_t(\lambda)\in\argmin_a\{c_t(a)+\lambda h_t(a)\}$.

\cue{Interpretation to say out loud.} $\lambda$ is a shadow price --- the
exchange rate between a unit of cost and a unit of error. Push $\lambda$ up and
error becomes ``expensive,'' so we defer more; push it down and we trust the AI.
A high-dimensional policy collapses to tuning one scalar.

\backup{Dual derivation if pressed: introduce $\lambda\ge0$ for the error
constraint and a \emph{free} multiplier $\mu_t\in\R$ for each simplex equality;
keep $x\ge0$ as the domain. Linearity in each $x_{t,a}$ forces the dual
feasibility condition $\mu_t\le\frac1T\{c_t(a)+\lambda h_t(a)\}$ for all $a$ (else
the inner inf is $-\infty$). The objective increases in each $\mu_t$, so
optimally $\mu_t^\star(\lambda)=\frac1T\min_a\{c_t(a)+\lambda h_t(a)\}$, which
substitutes to give $D(\lambda)$. Strong duality holds: the primal is feasible
(defer everything $\Rightarrow$ zero error) and bounded (costs $\ge0$), so
finite-dimensional LP duality applies. We assume losses in the argmin are
distinct so ties don't occur.}

% --------------------------------------------------------------------
\section{Threshold form}
\timecue{6:00}
\cue{Make the rule concrete --- this is the rule we actually run.} Under equal
cheap costs the chosen cheap model doesn't depend on $\lambda$, so the rule
factors into two steps. \textbf{Step 1, route:} pick the best cheap model
$j_t^\star\in\argmin_j\hath_t(j)$ and read off its score
$s_t=\min_j\hath_t(j)$ --- think of $s_t$ as an ``uncertainty.'' \textbf{Step 2,
threshold:} $\lambda$ only sets a cutoff $1/\lambda$:
\[
   s_t\ge 1/\lambda \;\Rightarrow\; \text{defer to expert},\qquad
   s_t< 1/\lambda \;\Rightarrow\; \text{accept the routed cheap label}.
\]

\cue{The payoff line.} With a single source and $\hath_t=$ uncertainty, this is
\emph{exactly} the original PAC-labeling threshold family. Multiple sources cost
nothing extra --- the score just becomes $\min_j\hath_t(j)$. And the very same
scalar will be \emph{fixed once} offline or \emph{moved over time} online ---
that's the fork in the talk.

\backup{We replace the unknown true loss by a proxy $\hath_t(a)$ --- an
uncertainty score, a calibrated error estimate, anything \emph{predictable} from
past audits and current non-expert info. Crucially proxies need \emph{not} be
calibrated or conservative: a better proxy improves cost, while held-out
calibration provides validity. In the experiments $\hath=1-\text{confidence}$.}

% --------------------------------------------------------------------
\section{Part I divider --- Offline}
\timecue{6:15}
\cue{One sentence.} ``In the offline case the whole pool is in front of us: we
spend a handful of expert labels to \emph{certify} one price, then label
everything at that fixed price.''

% --------------------------------------------------------------------
\section{Offline calibration}
\timecue{8:00}
\cue{The split that makes it work.} Decompose the price-$\lambda$ policy into its
realized loss $H_t(\lambda)=(1-D_t^\lambda)\,\ell(Y_t,\hatY_t^{(j_t^\star)})$ and
cost $C_t(\lambda)=D_t^\lambda$. The average \emph{cost} $C_T(\lambda)$ is fully
known from proxies and costs --- no labels needed. Only the average \emph{error}
$L_T(\lambda)=\frac1T\sum_t H_t(\lambda)$ needs true labels, so that is the one
thing we estimate. That is precisely what the calibration set buys.

\cue{Three steps.} (1) Sample $m$ calibration indices i.i.d.\ uniformly from the
pool (with replacement), pay for their expert labels, and form the empirical
loss $\hatL_m(\lambda)=\frac1m\sum_s H_{I_s}(\lambda)$. (2) Add a Hoeffding slack
for a one-sided upper confidence bound
$U_m(\lambda;\alpha)=\hatL_m(\lambda)+\sqrt{\log(1/\alpha)/2m}$. (3) Pick the
\emph{cheapest certified price}
$\hat\lambda=\min\{\lambda\ge0:U_m(\lambda;\alpha)\le\eps\}$, falling back to
expert-only ($\lambda=\infty$) if none qualifies.

\cue{Why we pay no union-bound tax --- this is the slick part.} As $\lambda$
increases, the threshold $1/\lambda$ \emph{decreases}, so each item can only move
from AI to expert, never between two cheap models. Hence $H_t(\lambda)$ is
non-increasing and $C_t(\lambda)$ non-decreasing in $\lambda$: the policies form
one \emph{nested, monotone} family. A single pointwise Hoeffding bound, placed at
the relevant crossing price, certifies the entire continuum --- no grid, no
$\log|\Lambda|$ penalty. Because $U_m$ is non-increasing and $C_T$
non-decreasing, the ``cheapest certified price'' is also the cost-minimizer.

\backup{Implementation: $m=0.2T$, $\alpha=0.05$, $\eps\in\{0.05,0.10,0.15\}$.
The sweep uses a grid (expert-only plus 201 evenly spaced thresholds) purely for
plotting --- the monotone certificate needs no grid. The original-PAC baseline
uses the same calibration sample and the same radius
$\sqrt{\log(1/\alpha)/2m}$ but picks its threshold from the calibration mistake
ordering rather than a grid.}

% --------------------------------------------------------------------
\section{Offline guarantee}
\timecue{9:00}
\cue{State the theorem plainly.} With probability $\ge1-\alpha$ over the
calibration draw, the realized average error on the \emph{whole} pool satisfies
$\frac1T\sum_t\ell(Y_t,\tildeY_t)\le\eps$. Finite-sample, distribution-free, and
it safely falls back to expert-only when no finite price qualifies.

\cue{Proof intuition in one breath.} Order the finitely many distinct policies by
$\lambda$. Let $\lambda^\circ$ be the last (fixed, data-independent) policy whose
\emph{true} loss still exceeds $\eps$. A bad selection $L_T(\hat\lambda)>\eps$
forces $\hat\lambda$ to land no later than $\lambda^\circ$; since $U_m$ is
non-increasing, $U_m(\lambda^\circ)\le U_m(\hat\lambda)\le\eps<L_T(\lambda^\circ)$.
So the bad event is contained in a \emph{single} fixed-price Hoeffding failure
$\{L_T(\lambda^\circ)>U_m(\lambda^\circ)\}$, which has probability $\le\alpha$.
Monotonicity is doing all the work --- that's why no union bound.

\cue{Generalization corollary --- mention briefly.} If calibration is i.i.d.\
from a distribution $P$, the \emph{population} risk at $\hat\lambda$ is $\le\eps$
w.p.\ $1-\alpha$; and on a fresh deployment pool of size $n$ the realized error
is $\le\eps+\sqrt{\log(1/\delta)/2n}$ w.p.\ $1-\alpha-\delta$. \cue{Repeat the
slogan:} proxies move cost, not validity --- a worse proxy just defers more
often.

% --------------------------------------------------------------------
\section{Offline experiment 1 --- one source}
\timecue{10:30}
\cue{Setup in a breath.} Two ImageNet prediction tables --- ImageNet
($T=50{,}000$) and ImageNetV2 ($T=10{,}000$) --- each with a ground-truth label,
a model prediction, and a confidence. Cheap source = the model prediction,
expert = ground truth, proxy loss $\hath=1-\text{confidence}$. Every number is
an average over \textbf{500 random trials}; the slide shows $\eps=0.10$.

\cue{Read the table left to right.} LP fixed-price vs.\ original PAC: error
$0.0873$ vs.\ $0.0879$ on ImageNet, budget saved $0.7848$ vs.\ $0.7864$ ---
essentially identical, with the \emph{violation column exactly zero}. That is the
sanity check: with one source the LP rule \emph{is} PAC, exactly as Slide 6
predicts. We skip the expert on $\sim$78\% of ImageNet items, $\sim$55\% on the
harder ImageNetV2.

\cue{Point at the figure.} Across $\eps\in\{0.05,0.10,0.15\}$ the two methods'
budget-saved and realized-error curves sit on top of each other; realized error
hugs --- and stays under --- the gray target line. Higher $\eps$ buys more
savings (0.63 $\to$ 0.79 $\to$ 0.89 on ImageNet).

% --------------------------------------------------------------------
\section{Offline experiment 2 --- two sources}
\timecue{11:30}
\cue{This is the slide that justifies the LP framing --- give it weight.} The
real data has only one genuine source, so we \emph{synthesize} two complementary
ones: partition items by the parity of the true class, and let each source reuse
the original predictor on its specialty half while emitting a random
low-confidence label off it. The pool is fixed across trials. Now per-item
routing genuinely matters.

\cue{The contrast --- the money numbers.} At $\eps=0.10$ on ImageNet, LP routing
saves \textbf{0.7848} of the expert budget versus \textbf{0.4746} for the
best single source --- and that baseline is already given a fair shake via a
Bonferroni correction (run both singles at $\alpha/2$, keep the better). Same
error control (violation $=0$), but the LP picks the right source per item and
defers far less. The pattern holds across $\eps$ and on ImageNetV2 (0.5457 vs.\
0.3587 at $\eps=0.10$).

\cue{Takeaway line.} ``A scalar threshold on a single source leaves a lot of
budget on the table; the LP recovers it by routing first, then thresholding.''

\backup{The monotone certificate still applies because, under equal costs, the
routed source $j_t^\star=\argmin_j\hath_t(j)$ doesn't depend on $\lambda$ --- the
score is just $\min_j\hath_t(j)$ and the policy is still nested in $\lambda$. The
grid here is 201 thresholds on $\min_j\hath_t(j)$ plus expert-only; same
$m=0.2T$, $\alpha=0.05$.}

% --------------------------------------------------------------------
\section{Part II divider --- Online}
\timecue{0:10 (segment-relative)}
\cue{One sentence.} ``Now items stream in one at a time, and an AI label's true
error is invisible unless we pay to audit it --- so we can no longer certify a
fixed price up front; the price has to \emph{move}.'' This begins your
$\approx$5-minute segment; the times below are relative to here.

% --------------------------------------------------------------------
\section{Online setup: moving threshold + audits}
\timecue{0:55}
\cue{The obstacle first.} The new difficulty versus offline: once we release an
AI label we normally \emph{never learn} whether it was right --- the loss
$\ell(Y_t,\hatY_t)$ is hidden unless we pay the expert. So we cannot just read
off our running error rate; we have to buy information.

\cue{Describe the loop concretely.} Each round hold a current price $\lambda_t$,
i.e.\ a threshold $1/\lambda_t$ on the routed uncertainty score
$s_t=\min_j\hath_t(j)$. Route to the lowest-proxy model; if it is too uncertain
($s_t\ge1/\lambda_t$) defer to the expert (always correct). Otherwise release the
AI label, then flip a coin: with probability $p$ \emph{audit} it --- pay the
expert, \emph{correct} the released label, and \emph{observe the AI's true loss}
$\ell_t$. Audits are rare on purpose, to stay cheap; that observed $\ell_t$ is
the only feedback, and it drives the next slide's update.

\backup{Predictability (Assumption~1): routing, the audit probability, and the
selected source use only past information plus the current non-expert info
$(X_t,\{\hatY_t^{(j)},\hath_t(j)\})$; given that, the audit draw
$A_t\sim\text{Bernoulli}(p_t)$ is independent, with $p_t\in[p_{\min},1]$ (in the
experiments $p_t\equiv p$ is constant). The proxy is also floored,
$\hath_t(j)\in[\rho,1]$ --- that floor is what bounds the price on Slide~15.}

% --------------------------------------------------------------------
\section{The price update}
\timecue{2:10}
\cue{The update is one-line dual ascent.}
$\lambda_{t+1}=[\lambda_t+\eta\widehat Z_t]_+$, projected to stay $\ge0$, where
$\widehat Z_t$ is this round's feedback. Three cases: $-\eps$ on a deferral;
$(\ell_t-\eps)/p$ on an audited AI label; $0$ on an unaudited AI label.

\cue{Narrate the sign.} An audited loss \emph{above} target ($\ell_t>\eps$) makes
$\widehat Z_t>0$, so $\lambda$ rises, the threshold $1/\lambda$ falls, and future
routing gets more cautious; a deferral or a clean audit ($\ell_t<\eps$) lowers
$\lambda$ and leans on the AI more.

\cue{Why the $1/p$ --- the key idea.} Audits are rare, so each audited item must
stand in for the $\sim\!1/p$ unaudited ones: inverse-probability weighting. That
makes the update \emph{unbiased} --- in expectation it behaves as if every AI
loss were seen, $\E[\widehat Z_t\mid\text{past}]=(\text{AI error})-\eps$.
\cue{Flag the simplification:} the update now centers \emph{directly} on the
final target $\eps$; there is no separate internal target $\eps_0$ any more (the
report dropped it). The only gap to $\eps$ is the vanishing slack on the next
slide.

\backup{In the proof's notation the three cases compress to the centered
inverse-observation signal $\widehat Z_t=\frac{O_t}{q_t}(H_t-\eps)$, with mask
$O_t=D_t+(1-D_t)A_t$ and $q_t=D_t+(1-D_t)p_t$ ($q_t=1$ on a deferral, $p_t$ on an
AI decision). Proposition~1: $\E[O_t/q_t\mid\calG_t]=1$ in both cases ($1/1$ and
$p_t/p_t$), hence $\E[\widehat Z_t\mid\calG_t]=H_t-\eps$. Keep $O_t,q_t,\calG_t$
off the slides --- they are bookkeeping for the martingale step only.}

% --------------------------------------------------------------------
\section{Finite-horizon guarantee: anatomy of the bound}
\timecue{3:35}
\cue{State the result.} For \emph{any} data sequence --- no distributional
assumptions, the only randomness is the audit coins --- with probability
$\ge1-\delta$,
\[
   \tfrac1T\sum_t\ell(Y_t,\tildeY_t)\;\le\;\eps+B_T,
   \qquad
   B_T=\frac{\Lambda_{\mathrm{on}}}{\eta T}
       +\frac{1}{p_{\min}}\sqrt{\frac{2\log(1/\delta)}{T}}.
\]
It controls \emph{realized} error, not regret against an offline oracle.

\cue{Read the bound term by term --- this is the slide's whole job.} (1) $\eps$
is the target the update steers toward. (2) $\Lambda_{\mathrm{on}}/(\eta T)$ is
optimization error --- how far the price can wander before it settles --- decaying
like $1/T$. (3) the last term is the \emph{price of partial feedback}: sampling
noise from auditing only a fraction, scaling like $1/\sqrt T$ and inflating as
the audit floor $p_{\min}$ shrinks.

\cue{Define $\Lambda_{\mathrm{on}}$ and give the intuition} (the new piece on the
slide). Explicitly $\Lambda_{\mathrm{on}}=\max\{\lambda_1,\,
1/\rho+\eta(1-\eps)/p_{\min}\}$ --- a constant ceiling the price can never exceed.
\emph{Why} it can't run away: the proxy is floored at $\rho$, so the moment
$\lambda_t$ reaches $1/\rho$ the threshold $1/\lambda_t$ drops to $\rho$ or below;
then \emph{every} item looks too uncertain and defers, which only pushes the
price back down. A single update can overshoot that level by at most one step
$\eta(1-\eps)/p_{\min}$ --- hence the formula. It is self-limiting.

\cue{The slack vanishes.} $B_T=O(T^{-1/2})$, so the released error approaches
$\eps$ as the stream grows; with $\eta\asymp T^{-1/2}$ that rate is explicit.
\cue{Released vs.\ AI error:} the bound is really on the counterfactual AI error,
but every audited mistake is corrected before release, so the shipped error is
only smaller.

\backup{Proof sketch (only if asked): (i) the price is bounded,
$0\le\lambda_t\le\Lambda_{\mathrm{on}}$, by the deferral argument above. (ii)
Since $[x]_+\ge x$, telescoping gives $\sum_t\widehat
Z_t\le\Lambda_{\mathrm{on}}/\eta$. (iii) The gap $M_t=(H_t-\eps)-\widehat Z_t$ is
a martingale-difference sequence (Prop.\ 1) bounded by $1/p_{\min}$, so
Azuma--Hoeffding gives $\sum_t M_t\le\frac{1}{p_{\min}}\sqrt{2T\log(1/\delta)}$.
(iv) Combine via $\sum_t H_t=T\eps+\sum_t\widehat Z_t+\sum_t M_t$, then use
audit-and-correct's pointwise $G_t\le H_t$ to pass from AI error to released
error.}

% --------------------------------------------------------------------
\section{Online experiment}
\timecue{4:20}
\cue{One clean message: $p$ is a validity--cost knob.} Per trial we randomly
permute the pool into a stream; settings $\eta=1$, proxy floor $\rho=0.01$,
$\lambda_1=1/\eps$, audit rate $p\in\{0.1,0.2,0.5\}$, again 500 trials. At
$\eps=0.10$ on ImageNet: $p=0.1$ saves $0.6629$ at final error $0.0780$;
$p=0.2$ saves $0.6115$ at $0.0737$; $p=0.5$ saves only $0.3952$ but drives final
error down to $0.0490$. Smaller $p$ = cheaper; larger $p$ = more corrections,
lower released error.

\cue{Point at the dashed vs.\ solid curves.} Dashed = counterfactual AI error
\emph{before} correction (hovers near $\eps$, this is the $H_t$ the theory
bounds); solid = realized error \emph{after} audit-and-correct (strictly below).
The gap between them is literally the value of auditing. The released-label error
stays controlled even under heavily partial feedback --- the online guarantee in
action.

\backup{Because the update centers on $\eps$ directly (no $\eps_0$ slack is
subtracted), the bound sits at exactly $\eps+B_T$, so at small $\eps$ and small
$p$ finite-$T$ fluctuation can produce a tiny violation (e.g.\ ImageNetV2,
$\eps=0.05$, $p=0.1$: violation $0.184$); it shrinks as $T$ grows or $p$ rises.
The ``AI error'' column is the pre-correction counterfactual; ``final error'' is
post-correction.}

% --------------------------------------------------------------------
\section{Takeaways}
\timecue{4:45}
\cue{Close on the unifying idea.} One LP dual gives a single price of error.
Offline we \emph{certify} it on held-out labels (Hoeffding + monotonicity);
online we \emph{move} it with audited inverse-probability feedback (dual ascent +
Azuma). In both regimes proxies affect efficiency, never validity.

\cue{The evidence in three beats.} LP matches PAC with one source; routing
clearly wins with two (0.78 vs.\ 0.47 budget saved); the online rule stays valid
under partial feedback with $p$ as the cost/error dial.

\cue{Honest limitation --- say it plainly.} Our real data has only one genuine
cheap source, so the two-source win uses a synthetic split. The genuinely
interesting regime is \emph{many sources with heterogeneous costs}: there the
routed source $j_t^\lambda=\argmin_j\{c_t(j)+\lambda\hath_t(j)\}$ depends on
$\lambda$, source-switching can make $L_T(\lambda)$ non-monotone, and the clean
single-point certificate breaks --- you fall back to certifying a finite price
grid (or enumerating breakpoints). Then thank the audience and invite questions.

% ====================================================================
\par\vspace{1.2\baselineskip}
\hrule
\vspace{0.4\baselineskip}
{\Large\bfseries\color{blue!45!black}Anticipated Questions \& Prepared Answers}
\vspace{0.3\baselineskip}

\qsection{Q1. Why bother with the LP at all if one source just recovers original PAC?}
Because the LP is what \emph{tells} you the right multi-source rule and where the
guarantee comes from. With one source it must coincide with PAC --- that's the
sanity check, not the contribution. The value shows up the moment you have more
than one source (or non-trivial costs): the dual hands you ``route to
$\argmin_a\{c_t(a)+\lambda\hath_t(a)\}$, then threshold'' automatically, and the
two-source experiment shows that buys $\sim$0.31 extra budget saved over the best
single source at the same error.

\qsection{Q2. The two cheap sources are synthetic --- is the gain real?}
Fair, and we flag it as the main limitation. The construction (specialty by
class parity, random off-specialty labels) is engineered so routing \emph{can}
help, so the number is an illustration of the mechanism, not a benchmark claim.
But the mechanism is real: whenever sources have complementary strengths, picking
the lower-proxy source per item before deferring strictly dominates committing to
one source. The honest next step is heterogeneous-cost data with several genuine
models.

\qsection{Q3. The expert is assumed perfect. What if experts are noisy?}
We treat the expert label as ground truth ($h_t(E)=0$) --- that's the standard
PAC-labeling abstraction. With noisy experts the ``error'' being certified is
error \emph{relative to the expert oracle}, and you'd need to model expert
disagreement (e.g.\ multiple experts, majority vote, or a known expert error
rate) to recover an absolute guarantee. Nothing in the LP forbids $h_t(E)>0$;
the clean fallback ``defer everything $\Rightarrow$ zero error'' is what would
need revisiting.

\qsection{Q4. Hoeffding is loose. Can you tighten the offline bound?}
Yes --- any valid one-sided upper confidence bound slots into the same monotone
argument. We note an empirical-Bernstein bound that uses the sample variance
$\widehat V_m(\lambda)$, which is much tighter when the calibration losses are
low-variance (which they are here, since most routed items are correct).
Betting-based confidence sequences also work if you reuse the calibration stream
while expanding the price set. We kept Hoeffding in the main results for
simplicity and comparability with PAC.

\qsection{Q5. Why no union bound over $\lambda$ --- isn't that cheating?}
No, it's the equal-cost structure. As $\lambda$ grows the threshold $1/\lambda$
only falls, so an item can only switch AI$\to$expert, never between cheap
sources. So $L_T(\lambda)$ is monotone and the candidate policies form one nested
chain; the bad-selection event collapses onto a single fixed crossing price, so
one pointwise Hoeffding bound suffices. The moment costs are unequal, source
selection depends on $\lambda$, monotonicity breaks, and you \emph{do} pay a
union bound over a price grid (or over the $\le1+(T+m)\binom{k+1}{2}$ breakpoints).

\qsection{Q6. Why does the online error sometimes exceed $\eps$?}
The guarantee is $\frac1T\sum_t\ell(Y_t,\tildeY_t)\le\eps+B_T$, with
$B_T=\Lambda_{\mathrm{on}}/(\eta T)+\tfrac1{p_{\min}}\sqrt{2\log(1/\delta)/T}$.
The price update centers on $\eps$ \emph{directly}, so the bound sits exactly at
$\eps+B_T$, not at $\eps$: there is a finite-horizon cushion $B_T$ that the
target is allowed to overshoot by. At small $\eps$ and small $p$ (few audits,
larger $B_T$) finite-$T$ noise can push the realized error a little above $\eps$
--- e.g.\ ImageNetV2, $\eps=0.05$, $p=0.1$. Since $B_T=O(T^{-1/2})$, the cushion
shrinks as the stream lengthens or the audit rate rises, and the error settles
back to $\eps$. (An earlier draft carried a separate internal target $\eps_0$ and
``spent'' this slack explicitly; the current version folds it into $B_T$.)

\qsection{Q7. How do you choose the audit rate $p$ (and the step size $\eta$)?}
$p$ is a deliberate cost/validity dial: small $p$ saves budget but gives noisier
feedback (larger $1/p_{\min}$ term); large $p$ corrects more labels and lowers
realized error. The floor $p_{\min}$ enters the bound as $1/p_{\min}$, so you
don't want it too small. For $\eta$, the bound is
$\Lambda_{\mathrm{on}}/(\eta T)+O(1/\sqrt T)$, optimized at
$\eta\asymp T^{-1/2}$, which makes the whole slack $O(T^{-1/2})$. In the runs we
used a fixed $\eta=1$ and still saw control because $T$ is large.

\qsection{Q8. If the dashed counterfactual AI-error curve already satisfies the error bound, are audits actually needed?}
The dashed curve is not ``use AI on every item.'' It is the error of the
\emph{selected} AI labels after the deferral rule has already removed
high-proxy-loss items: $H_t=(1-D_t)\ell(Y_t,\hat Y_t)$. So the fact that it is
below $\eps$ says the confidence proxy and threshold/router are working well in
this dataset. It does \emph{not} mean the online learner could certify or adapt
without labels. Audits provide the feedback for updating $\lambda_t$ and protect
against miscalibration or stream shift; without audits, AI-label losses are
hidden and a bad threshold could silently violate the target. The right takeaway
is that a strong proxy may reduce the number of audits needed, not that audits
are unnecessary.

\qsection{Q9. Is this regret or error control? How does it relate to online conformal prediction?}
It's \emph{realized-error} control over a fixed horizon, not regret against an
offline expert-query oracle. The lineage is online/adaptive conformal prediction
with partial (semi-bandit) feedback --- those methods update a calibration
parameter online, and the key borrowed idea is that when labels aren't always
revealed you must correct the update for the observation process. Our
inverse-probability signal $O_t/q_t$ is exactly that correction. The object we
control is the average error of the \emph{released dataset}, not prediction-set
coverage.

\qsection{Q10. What's the actual cost of calibration / auditing --- isn't that hidden expert spend?}
Good --- we separate three costs in the report. \emph{Deployment} cost is
$C_T(\lambda)$, the quantity we minimize. \emph{Development} cost is the expert
labels used to train the proxy and to certify the price ($m=0.2T$ here); if those
labeled items are part of the pool and corrected in the release, there's no
duplicate expert charge for them. \emph{Audit} cost online is
$\E[C_t\mid\calG_t]=D_t+(1-D_t)p_t$ per round --- the extra expert spend to obtain
feedback. So the headline ``78\% saved'' is deployment savings; calibration is a
one-time, accountable overhead.

\qsection{Q11. Does the guarantee need i.i.d.\ data?}
The offline pool guarantee is \emph{conditional on the fixed pool} --- it only
needs the calibration indices to be sampled i.i.d.\ from the pool, not the data
to be i.i.d.\ The \emph{generalization} corollary (population risk + fresh
deployment pool) does assume an i.i.d.\ source distribution. The online guarantee
is fully \emph{distribution-free in the data sequence} --- it can even be
adversarial --- with all randomness coming from the audit coin flips. That
robustness is one of the nicer features of the online side.

\qsection{Q12. Why is ImageNetV2 harder (lower savings, occasional violation)?}
ImageNetV2 is a known distribution shift from ImageNet --- the base model is less
accurate and less calibrated on it, so at the same $\eps$ fewer items clear the
confidence threshold and we defer more (savings drop from $\sim$0.78 to
$\sim$0.55 at $\eps=0.10$). It's a useful stress test precisely because the proxy
is worse there, and the guarantee still holds --- consistent with ``proxies
affect cost, not validity.''

\end{document}
